\documentclass[10pt,a4paper]{amsart}
\usepackage{amsmath,amssymb,amsthm,graphicx}
% Retained preamble for the Stacks Project chapter snapshot in this directory.
% The chapter's own preamble.tex is a copy of the shared Stacks Project preamble; the block
% below is sliced MECHANICALLY, line for line, from it. Nothing was added, removed or
% redefined; the wrapper class, the xr-hyper/hyperref wiring and the externaldocument list
% are omitted because the assembler supplies the document class.
%
% lines taken from preamble.tex: 3, 5, 9, 11, 12, 156, 157, 158, ... (51 lines total)
%
% package coverage restored from the chapter's own preamble (2026-10-13): xy, xy-2cell

\usepackage[all]{xy}
\xyoption{2cell}
\UseAllTwocells

\newenvironment{reference}{\comment}{\endcomment}
\newenvironment{slogan}{\comment}{\endcomment}
\newenvironment{history}{\comment}{\endcomment}
\theoremstyle{plain}
\newtheorem{theorem}[subsection]{Theorem}
\newtheorem{proposition}[subsection]{Proposition}
\newtheorem{lemma}[subsection]{Lemma}
\theoremstyle{definition}
\newtheorem{definition}[subsection]{Definition}
\newtheorem{example}[subsection]{Example}
\newtheorem{exercise}[subsection]{Exercise}
\newtheorem{situation}[subsection]{Situation}
\theoremstyle{remark}
\newtheorem{remark}[subsection]{Remark}
\newtheorem{remarks}[subsection]{Remarks}
\numberwithin{equation}{subsection}
\def\lim{\mathop{\mathrm{lim}}\nolimits}
\def\colim{\mathop{\mathrm{colim}}\nolimits}
\def\Spec{\mathop{\mathrm{Spec}}}
\def\Hom{\mathop{\mathrm{Hom}}\nolimits}
\def\Ext{\mathop{\mathrm{Ext}}\nolimits}
\def\SheafHom{\mathop{\mathcal{H}\!\mathit{om}}\nolimits}
\def\SheafExt{\mathop{\mathcal{E}\!\mathit{xt}}\nolimits}
\def\Sch{\mathit{Sch}}
\def\Mor{\mathop{\mathrm{Mor}}\nolimits}
\def\Ob{\mathop{\mathrm{Ob}}\nolimits}
\def\Sh{\mathop{\mathit{Sh}}\nolimits}
\def\NL{\mathop{N\!L}\nolimits}
\def\CH{\mathop{\mathrm{CH}}\nolimits}
\def\proetale{{pro\text{-}\acute{e}tale}}
\def\etale{{\acute{e}tale}}
\def\QCoh{\mathit{QCoh}}
\def\Ker{\mathop{\mathrm{Ker}}}
\def\Im{\mathop{\mathrm{Im}}}
\def\Coker{\mathop{\mathrm{Coker}}}
\def\Coim{\mathop{\mathrm{Coim}}}
\DeclareMathSymbol{\boxtimes}{\mathbin}{AMSa}{"02}
\def\QCohstack{\mathcal{QC}\!\mathit{oh}}
\def\Cohstack{\mathcal{C}\!\mathit{oh}}
\def\Spacesstack{\mathcal{S}\!\mathit{paces}}
\def\Quotfunctor{\mathrm{Quot}}
\def\Hilbfunctor{\mathrm{Hilb}}
\def\Curvesstack{\mathcal{C}\!\mathit{urves}}
\def\Polarizedstack{\mathcal{P}\!\mathit{olarized}}
\def\Complexesstack{\mathcal{C}\!\mathit{omplexes}}
\def\Pic{\mathop{\mathrm{Pic}}\nolimits}
\def\Picardstack{\mathcal{P}\!\mathit{ic}}
\def\Picardfunctor{\mathrm{Pic}}
\def\Deformationcategory{\mathcal{D}\!\mathit{ef}}
\renewenvironment{reference}{}{}
\renewenvironment{slogan}{}{}
\renewenvironment{history}{}{}
\usepackage{cleveref}
\crefname{originaltheorem}{original theorem}{original theorems}
\Crefname{originaltheorem}{Original theorem}{Original theorems}
\makeatletter
\newlabel{lemma-algebraize-local-cohomology}{{0DXP}{source}{}{sourcepointer-1}{}}
\newlabel{lemma-algebraize-local-cohomology@cref}{{[originaltheorem][0][]0DXP}{[1][0][]source}{}{}{}}
\makeatother
\crefname{originaltheorem}{original theorem}{original theorems}
\Crefname{originaltheorem}{Original theorem}{Original theorems}
\makeatletter
\newlabel{lemma-compare-with-derived-completion}{{0DXI}{source}{}{sourcepointer-2}{}}
\newlabel{lemma-compare-with-derived-completion@cref}{{[originaltheorem][0][]0DXI}{[1][0][]source}{}{}{}}
\makeatother
\crefname{originaltheorem}{original theorem}{original theorems}
\Crefname{originaltheorem}{Original theorem}{Original theorems}
\makeatletter
\newlabel{lemma-kill-colimit}{{0DXN}{source}{}{sourcepointer-3}{}}
\newlabel{lemma-kill-colimit@cref}{{[originaltheorem][0][]0DXN}{[1][0][]source}{}{}{}}
\makeatother
\crefname{originaltheorem}{original theorem}{original theorems}
\Crefname{originaltheorem}{Original theorem}{Original theorems}
\makeatletter
\newlabel{local-cohomology-lemma-finiteness-pushforwards-and-H1-local}{{0BK0}{source}{}{sourcepointer-4}{}}
\newlabel{local-cohomology-lemma-finiteness-pushforwards-and-H1-local@cref}{{[originaltheorem][0][]0BK0}{[1][0][]source}{}{}{}}
\makeatother
\title{Formal-sections algebraization theorem: cohomology of F on small opens equals the formal limit below degree s, with finiteness under a stronger depth bound}
\author{The Stacks Project Authors}
\date{stacks-project:algebraization.tex:local-snapshot-20261012; GNU FDL 1.2 or later}
\begin{document}
\maketitle
\noindent\textit{Editorial note.} Statement span is lines 3768-3808 of sources/stacks-algebraization/pinned-original-source.tex; proof span is lines 3810-3846, the complete author proof (35 body lines).\par\medskip
\noindent\textit{Editorial note.} Navigation stubs: 4 cited-but-not-included labels are supplied as non-load-bearing pointers with their pinned tags where the snapshot tag index has them.\par\medskip
\noindent\textit{Editorial note.} extra\_packages is empty on purpose: the retained preamble already loads xy with options and a second declaration clashes in the assembler wrapper.\par\medskip
\noindent\textit{Editorial note.} Text is retained verbatim, including the author's own typos and placeholders.\par\medskip
\noindent\textit{Editorial note.} Pool algebraization wave36 (al02 bib-salvage batch). Spans computed from \textbackslash{}label (16.5/16.6).\par\medskip
\noindent\textit{Editorial note.} External \textbackslash{}cite keys salvaged with card-local bibliography.tex entries copied verbatim from the snapshot my.bib (precedent S3-morph-blow-up-ample-family); URL/ISSN/EPRINT fields dropped as the preamble has no url package.\par\medskip
\noindent\textit{Editorial note.} statement contains reference environment with real content -> patch-sd-envs.tex applied (sd02 lesson).\par\medskip
\noindent\textit{Allowed standard prerequisites (editorial).}\begin{itemize}
\item Explicit prerequisite boundary. See why-field; in-chapter lemmas cited by name become stub prereqs; external cites resolved via the card-local bibliography.tex fragment.
\end{itemize}\medskip
\ifdefined\hypertarget\hypertarget{sourcepointer-1}{}\fi
\noindent\textit{Source reference.} Original theorem 0DXP: sources/stacks-algebraization/pinned-original-source.tex, source lines 2720--2735. Cited by the retained statement or proof but not part of the collected spans; supplied only as a navigation pointer. The argument does not depend on its content.\par\medskip
\ifdefined\hypertarget\hypertarget{sourcepointer-2}{}\fi
\noindent\textit{Source reference.} Original theorem 0DXI: sources/stacks-algebraization/pinned-original-source.tex, source lines 3736--3746. Cited by the retained statement or proof but not part of the collected spans; supplied only as a navigation pointer. The argument does not depend on its content.\par\medskip
\ifdefined\hypertarget\hypertarget{sourcepointer-3}{}\fi
\noindent\textit{Source reference.} Original theorem 0DXN: sources/stacks-algebraization/pinned-original-source.tex, source lines 2688--2694. Cited by the retained statement or proof but not part of the collected spans; supplied only as a navigation pointer. The argument does not depend on its content.\par\medskip
\ifdefined\hypertarget\hypertarget{sourcepointer-4}{}\fi
\noindent\textit{Source reference.} Original theorem 0BK0: sources/stacks-algebraization/pinned-original-source.tex, source lines 3846--3846. Cited by the retained statement or proof but not part of the collected spans; supplied only as a navigation pointer. The argument does not depend on its content.\par\medskip
\begin{theorem}
\label{theorem-algebraization-formal-sections}
\begin{reference}
The method of proof follows roughly the method of
proof of \cite[Theorem 1]{Faltings-algebraisation}
and \cite[Satz 2]{Faltings-uber}.
The result is almost the same as
\cite[Theorem 1.1]{MRaynaud-paper} (affine complement case) and
\cite[Theorem 3.9]{MRaynaud-book} (complement is union of few affines).
\end{reference}
Let $(A, \mathfrak m)$ be a Noetherian local ring which has a
dualizing complex and is complete with respect to an ideal $I$.
Set $X = \Spec(A)$, $Y = V(I)$, and $U = X \setminus \{\mathfrak m\}$.
Let $\mathcal{F}$ be a coherent sheaf on $U$.
Assume
\begin{enumerate}
\item $\text{cd}(A, I) \leq d$, i.e.,
$H^i(X \setminus Y, \mathcal{G}) = 0$ for $i \geq d$ and
quasi-coherent $\mathcal{G}$ on $X$,
\item for any $x \in X \setminus Y$ whose closure $\overline{\{x\}}$
in $X$ meets $U \cap Y$ we have
$$
\text{depth}_{\mathcal{O}_{X, x}}(\mathcal{F}_x) \geq s
\quad\text{or}\quad
\text{depth}_{\mathcal{O}_{X, x}}(\mathcal{F}_x)
+ \dim(\overline{\{x\}}) > d + s
$$
\end{enumerate}
Then there exists an open $V_0 \subset U$ containing $U \cap Y$
such that for any open $V \subset V_0$ containing $U \cap Y$
the map
$$
H^i(V, \mathcal{F}) \to \lim H^i(U, \mathcal{F}/I^n\mathcal{F})
$$
is an isomorphism for $i < s$. If in addition
$
\text{depth}_{\mathcal{O}_{X, x}}(\mathcal{F}_x) +
\dim(\overline{\{x\}}) > s
$
for all $x \in U \cap Y$, then these cohomology groups are finite $A$-modules.
\end{theorem}

\begin{proof}
Choose a finite $A$-module $M$ such that $\mathcal{F}$ is the
restriction to $U$ of the
coherent $\mathcal{O}_X$-module associated to $M$, see Local Cohomology,
Lemma \ref{local-cohomology-lemma-finiteness-pushforwards-and-H1-local}.
Then the assumptions of
Lemma \ref{lemma-algebraize-local-cohomology}
are satisfied.
Pick $J_0$ as in that lemma and set $V_0 = X \setminus V(J_0)$.
Then opens $V \subset V_0$ containing $U \cap Y$
correspond $1$-to-$1$ with ideals $J \subset J_0$ with
$V(J) \cap V(I) = \{\mathfrak m\}$.
Moreover, for such a choice we have a distinguished triangle
$$
R\Gamma_J(M) \to M \to R\Gamma(V, \mathcal{F}) \to
R\Gamma_J(M)[1]
$$
We similarly have a distinguished triangle
$$
R\Gamma_\mathfrak m(M)^\wedge \to
M \to
R\Gamma(U, \mathcal{F})^\wedge \to
R\Gamma_\mathfrak m(M)^\wedge[1]
$$
involving derived $I$-adic completions.
The cohomology groups of $R\Gamma(U, \mathcal{F})^\wedge$ are
equal to the limits in the statement of the theorem by
Lemma \ref{lemma-compare-with-derived-completion}.
The canonical map between these triangles
and some easy arguments show that our
theorem follows from the main Lemma \ref{lemma-algebraize-local-cohomology}
(note that we have $i < s$ here whereas we have
$i \leq s$ in the lemma; this is because of the shift).
The finiteness of the cohomology groups
(under the additional assumption) follows from
Lemma \ref{lemma-kill-colimit}.
\end{proof}

\begin{thebibliography}{9}
\bibitem{Faltings-algebraisation}
Faltings, Gerd.
\newblock \emph{Algebraisation of some formal vector bundles}.
\newblock Ann. of Math. (2) 110 (3), 501--514 (1979).
\bibitem{Faltings-uber}
Faltings, Gerd.
\newblock \emph{\"Uber lokale {K}ohomologiegruppen hoher {O}rdnung}.
\newblock J. Reine Angew. Math. 313, 43--51 (1980).
\bibitem{MRaynaud-book}
Raynaud, Mich\`ele.
\newblock \emph{Th\'eor\`emes de {L}efschetz en cohomologie coh\'erente et en cohomologie \'etale}.
\newblock Soci\'et\'e Math\'ematique de France, Paris, 1975.
\newblock Bull. Soc. Math. France, M\'em. No. 41, Suppl\'ement au Bull. Soc. Math. France, Tome 103.
\bibitem{MRaynaud-paper}
Raynaud, Mich\`ele.
\newblock \emph{Th\'eoremes de {L}efschetz en cohomologie des faisceaux coh\'erents et en cohomologie \'etale. {A}pplication au groupe fondamental}.
\newblock Ann. Sci. \'Ecole Norm. Sup. (4) 7, 29--52 (1974).
\newblock Collection of articles dedicated to Henri Cartan on the occasion of his 70th birthday, I.
\end{thebibliography}
\end{document}
